\documentclass[10pt,a4paper]{amsart}
\usepackage[margin=19mm]{geometry}
\usepackage{amsmath,amssymb,mathtools}
\usepackage[scr=rsfs,cal=euler]{mathalfa}
\usepackage{xfrac}
\usepackage[unicode,hidelinks,bookmarks=false]{hyperref}
\newcommand{\ie}{i.e.\ }
\newcommand{\cf}{cf.\ }
\newcommand{\resp}{respectively\ }
\newcommand{\Subsection}{\subsection}
\providecommand{\nobreakdash}{\nobreak\hbox{-}}
\setlength{\emergencystretch}{2em}
\multlinegap0pt
\usepackage{bm}
\usepackage{bbm}
\usepackage{dsfont}
\usepackage{xspace}
\usepackage{cancel}
\usepackage{mathtools}
\usepackage{amsthm}
\makeatletter
\@ifundefined{lem}{\newtheorem{lem}{Lemma}}{}
\@ifundefined{thm}{\newtheorem{thm}{Theorem}}{}
\makeatother
\title[Eigenvalues and eigenfunctions of the non-local dispersal wi]{Eigenvalues and eigenfunctions of the non-local dispersal with Neumann-type boundary condition and symmetric kernel}
\author{Maciej Tadej}
\date{Source: arXiv:2605.25154v1 (CC BY 4.0)}
\begin{document}
\maketitle

\noindent\emph{Source and licence.} Eigenvalues and eigenfunctions of the non-local dispersal with Neumann-type boundary condition and symmetric kernel. Version arXiv:2605.25154v1 (CC BY 4.0), distributed under the Creative Commons Attribution 4.0 International licence, as recorded in the arXiv licence metadata for this exact version and shown on the abstract page https://arxiv.org/abs/2605.25154v1. Full rights notice: licenses/arxiv-2605.25154-CC-BY-4.0.txt.\par\smallskip

\noindent\textit{Editorial scope.} Counted unit: the Lem whose source label is \texttt{lem:convergence\_principal\_eigenfunction} in the article (source0.tex lines 691--697, source label \texttt{lem:convergence\_principal\_eigenfunction}), delivered together with the article's own complete original proof of it (source0.tex lines 699--798, one \texttt{proof} environment of 100 lines). No mathematics is written, repaired or re-derived: statement and proof are the article's own text line for line. Required context retained uncounted: thm:contninous\_spectrum (lines 272--278); eq:v\_N\_Def (lines 622--625); eq:variational\_discrete (lines 654--657); lem:convergence\_principal\_eigenvalue (lines 669--675). Every definition, display and label of those lines is retained unchanged. Omitted from this package and not redistributed: 1--271, 279--621, 626--653, 658--668, 676--690, 698--698, 799--878. Nothing inside a retained excerpt is omitted or altered: every retained body is delivered exactly as the article's lines are written, so no omission receipt of any kind accompanies this package. Citations: the retained excerpts carry no citation command at all, so no bibliography entry of the article is transcribed and no reference is redistributed. Reference adaptation: none was needed and none was made; every reference inside the delivered unit resolves inside it, so no unresolved reference or citation is produced. Printed slips kept literal: no printed word, symbol, hypothesis or number of the counted statement or of its proof was altered. Layout and class: the article's own class is \texttt{article} with packages including geometry, subcaption, graphicx, float, titlesec, amsmath, amssymb, amsthm, cases, babel, hyperref, enumitem, authblk, dsfont; the wrapper is the lane's pinned \texttt{amsart} editorial wrapper at 10pt on A4, which restates the theorem-like environments the retained text needs and adds the article's own macro definitions that the retained text needs (none), with extra wrapper packages (bm, bbm, dsfont, xspace, cancel, mathtools). The delivered numbering is the unit's own. Assets: no retained span contains a figure environment, a graphics-inclusion command, a PDF-page inclusion, a table or a TikZ picture; no image file is copied into this package, the package declares no asset dependency and no image file is redistributed with it. Licence: version arXiv:2605.25154v1 of this article is distributed under the Creative Commons Attribution 4.0 International licence, as recorded in the arXiv licence metadata for this exact version and shown on the abstract page https://arxiv.org/abs/2605.25154v1. A full rights notice is delivered at licenses/arxiv-2605.25154-CC-BY-4.0.txt. Characters: every retained excerpt is pure ASCII, so the delivered engine, XeLaTeX with its default Latin Modern text font, reports no missing character.
\begin{thm}
    \label{thm:contninous_spectrum}
    Let $\Omega$ and $J$ satisfy Assumption 1.1. The essential (continuous) spectrum of the non-local operator $\mathcal{L}$, denoted by $\sigma_c := \sigma_c(\mathcal{L})$, is given by the range of its multiplication part
    \begin{equation}
        \sigma_c = \left\{ -\int_\Omega J(x-y) \,dy : x \in \overline{\Omega} \right\}.
    \end{equation}
\end{thm}

    \begin{equation}
        \label{eq:v_N_Def}
        v_1^N(x) = \sum_{m=0}^N c_m \varphi_m(x).
    \end{equation}

    \begin{equation}
        \label{eq:variational_discrete}
        \beta_1^N := \sup_{\mathbf{c} \in \mathcal{C}} \: \langle (\mathbf{A} - \mathbf{B}) \mathbf{c}, \mathbf{c} \rangle
    \end{equation}

\begin{lem}
    \label{lem:convergence_principal_eigenvalue}
    Let $\beta_1^N$ be the principal eigenvalue of the discrete problem \eqref{eq:variational_discrete}. Then
    \begin{equation*}
        \beta_1^N \rightarrow \beta_1 \quad \text{ as } N \rightarrow \infty.
    \end{equation*}
\end{lem}

\begin{lem}
    \label{lem:convergence_principal_eigenfunction}
    Let $v_1^N$ be the function defined by formula \eqref{eq:v_N_Def} with vector of coefficients $\mathbf{c}$ being the maximizer of the problem \eqref{eq:variational_discrete}. Then 
    \begin{equation*}
        v_1^N \rightarrow v_1 \quad \text{ as } N \rightarrow \infty.
    \end{equation*}
\end{lem}

\begin{proof}
    Recall that $\|v_1^N\|_{L^2} = 1$. By the Banach-Alaoglu theorem, there exists a subsequence, still denoted by $v_1^N$, and a function $v \in L^2(\Omega)$ such that $v_1^N \rightharpoonup v_1$ weakly in $L^2(\Omega)$. The zero mass constraint is preserved under weak convergence
    \begin{equation}
        \int_\Omega v_1(x) \:dx = \lim_{N \to \infty} \int_\Omega v_1^N(x) \:dx = 0.
    \end{equation}
    To identify the limit, consider an arbitrary test function $\psi \in V_M$. For $N \geq M$, the discrete equation reads $\langle \mathcal{L} v_1^N, \psi \rangle = \beta_1^N \langle v_1^N, \psi \rangle$. Passing to the limit $N \to \infty$, and using the self-adjointness of $\mathcal{L}$, we obtain
    \begin{align}
        \langle v_1, \mathcal{L} \psi \rangle &= \lim_{N \to \infty} \langle v_1^N, \mathcal{L} \psi \rangle = \lim_{N \to \infty} \langle \mathcal{L} v_1^N, \psi \rangle \\
        &= \lim_{N \to \infty} \beta_1^N \langle v_1^N, \psi \rangle = \beta_1 \langle v_1, \psi \rangle.
    \end{align}
    Since $\bigcup_{M} V_M$ is dense in $L^2(\Omega)$, this implies $\mathcal{L} v_1 = \beta_1 v_1$ in the weak sense. It remains to prove that the weak convergence $v_1^N \rightharpoonup v_1$ is strong in $L^2(\Omega)$. We decompose the non-local operator $\mathcal{L}$ into a compact integral part $\mathcal{K}$ and a local multiplication part $\mathcal{M}$:
    \begin{equation}
        \mathcal{L}u(x) = (\mathcal{K}u)(x) - (\mathcal{M}u)(x) = \int_\Omega J(x-y)u(y)\,dy - b(x)u(x),
    \end{equation}
    where $b(x) = \int_\Omega J(x-y)dy$. Theorem \ref{thm:contninous_spectrum} states that the essential spectrum is generated by the multiplication operator, yielding $\sup \sigma_c = \max_{x \in \overline{\Omega}} (-b(x))$. From the Galerkin approximation, we evaluate the energy sequence
    \begin{equation}
        \langle \mathcal{L} v_1^N, v_1^N \rangle = \beta_1^N \to \beta_1 \quad \text{as } N \to \infty.
    \end{equation}
    Using our decomposition of $\mathcal{L}$, we can rewrite this energy equation as:
    \begin{equation}
        \label{eq:weak_solution_sequence}
        \langle \mathcal{K} v_1^N, v_1^N \rangle - \int_\Omega b(x) (v_1^N(x))^2 \, dx = \beta_1^N.
    \end{equation}
    Rearranging the terms in equation \eqref{eq:weak_solution_sequence} to group the $L^2$ norms, we get
    \begin{equation} \label{eq:rearranged_energy}
        \int_\Omega (\beta_1^N + b(x)) (v_1^N(x))^2 \, dx = \langle \mathcal{K} v_1^N, v_1^N \rangle.
    \end{equation}
    
    Because $J \in C^0(\mathbb{R}^n)$ and $\Omega$ is bounded, the integral operator $\mathcal{K} : L^2(\Omega) \to L^2(\Omega)$ is a compact operator. A fundamental property of compact operators is that they map weakly convergent sequences into strongly convergent ones. Since $v_1^N \rightharpoonup v_1$ weakly in $L^2(\Omega)$, it follows that $\mathcal{K} v_1^N \to \mathcal{K} v_1$ strongly. This guarantees the convergence of the inner product
    \begin{equation}
        \label{eq:scalar_product_compact_operator}
        \lim_{N \to \infty} \langle \mathcal{K} v_1^N, v_1^N \rangle = \langle \mathcal{K} v_1, v_1 \rangle.
    \end{equation}
    Next, from Lemma \ref{lem:convergence_principal_eigenvalue} we have that $\lim_{N \to \infty} \beta_1^N = \beta_1$ thus we write
    \begin{align}
        \lim_{N \to \infty} \int_\Omega (\beta_1^N + b(x)) (v_1^N(x))^2 \, dx 
        &= \lim_{N \to \infty} \left( \beta_1^N \int_\Omega (v_1^N(x))^2 \, dx + \int_\Omega b(x) (v_1^N(x))^2 \, dx \right) \nonumber \\
        &= \lim_{N \to \infty} \left( \beta_1^N \cdot 1 + \int_\Omega b(x) (v_1^N(x))^2 \, dx \right) \nonumber \\
        &= \beta_1 + \lim_{N \to \infty} \int_\Omega b(x) (v_1^N(x))^2 \, dx \nonumber \\
        &= \lim_{N \to \infty} \left( \beta_1 \int_\Omega (v_1^N(x))^2 \, dx + \int_\Omega b(x) (v_1^N(x))^2 \, dx \right) \nonumber \\
        &= \lim_{N \to \infty} \int_\Omega (\beta_1 + b(x)) (v_1^N(x))^2 \, dx.
        \label{eq:weighted_L2_norm}
    \end{align}
    Applying formulas \eqref{eq:scalar_product_compact_operator} and \eqref{eq:weighted_L2_norm} into equation \eqref{eq:rearranged_energy} we obtain
    \begin{equation} 
        \label{eq:limit_compact}
        \int_\Omega (\beta_1 + b(x)) v_1^N(x)^2 \, dx = \langle \mathcal{K} v_1, v_1 \rangle.
    \end{equation}
    Now, recall that the weak limit $v_1$ satisfies the equation $\mathcal{L} v_1 = \beta_1 v_1$. By rewriting this exact relation using the operator decomposition and taking the inner product with $v_1$, we obtain:
    \begin{equation}
        \langle \mathcal{K} v_1, v_1 \rangle - \int_\Omega b(x) v_1(x)^2 \, dx = \beta_1 \int_\Omega v_1(x)^2 \, dx.
    \end{equation}
    Rearranging this yields
    \begin{equation} \label{eq:weak_limit_energy}
        \int_\Omega (\beta_1 + b(x)) v_1(x)^2 \, dx = \langle \mathcal{K} v_1, v_1 \rangle.
    \end{equation}
    Combining equations \eqref{eq:limit_compact} and \eqref{eq:weak_limit_energy}, we have rigorously proven that
    \begin{equation} \label{eq:norm_convergence}
        \lim_{N \to \infty} \int_\Omega (\beta_1 + b(x)) (v_1^N(x))^2 \, dx = \int_\Omega (\beta_1 + b(x)) v_1(x)^2 \, dx.
    \end{equation}
    
    Finally, we utilize the crucial spectral gap condition. From Proposition 2.1, we established that $\beta_1 > \sup \sigma_c = \max_{x \in \overline{\Omega}} (-b(x))$. This implies that for all $x \in \overline{\Omega}$:
    \begin{equation}
        \beta_1 + b(x) \geq \beta_1 - \sup \sigma_c > 0.
    \end{equation}
    Because the weight function $w(x) := \beta_1 + b(x)$ is strictly positive and bounded, it induces an equivalent norm on the Hilbert space $L^2(\Omega)$, defined by
    \begin{equation}
        \|u\|_{*} := \left( \int_\Omega (\beta_1 + b(x)) u(x)^2 \, dx \right)^{1/2}.
    \end{equation}
    
    Equation \eqref{eq:norm_convergence} demonstrates exactly that the integral of the squared sequence with this weight converges to the integral of the squared limit. 
    
    To rigorously prove that $v_1^N$ converges strongly to $v_1$ in the standard $L^2(\Omega)$ norm, we evaluate the limit of the weighted squared difference $v_1^N - v_1$. Because $\beta_1 + b(x) \geq c > 0$, we can bound the standard $L^2$ norm by this weighted integral. Expanding the square inside the integral yields:
    \begin{align}
        0 \leq c \|v_1^N - v_1\|_{L^2}^2 &= c \int_\Omega (v_1^N(x) - v_1(x))^2 \, dx \nonumber \\
        &\leq \int_\Omega (\beta_1 + b(x))(v_1^N(x) - v_1(x))^2 \, dx \nonumber \\
        &= \int_\Omega (\beta_1 + b(x))(v_1^N(x))^2 \, dx \nonumber \\
        &\quad - 2 \int_\Omega (\beta_1 + b(x))v_1^N(x)v_1(x) \, dx \nonumber \\
        &\quad + \int_\Omega (\beta_1 + b(x))(v_1(x))^2 \, dx.
    \end{align}
    
    We now take the limit as $N \to \infty$ for each of the three terms on the right-hand side. The limit of the first was already established in equation \eqref{eq:norm_convergence}. For the second term, notice that since the weight function is bounded  and $v_1^N \rightharpoonup v_1 \in L^2(\Omega)$, we have that
    \begin{equation}
        \lim_{N \to \infty} \int_\Omega v_1^N(x) \big[ (\beta_1 + b(x)) v_1(x) \big] \, dx = \int_\Omega v_1(x) \big[ (\beta_1 + b(x)) v_1(x) \big] \, dx.
    \end{equation}
    The third term is independent of $N$. Combining these limits yields that
    \begin{align}
        \lim_{N \to \infty} \int_\Omega (\beta_1 + b(x))(v_1^N(x) - v_1(x))^2 \, dx 
        &= \int_\Omega (\beta_1 + b(x))(v_1(x))^2 \, dx \nonumber \\
        &\quad - 2 \int_\Omega (\beta_1 + b(x))(v_1(x))^2 \, dx \nonumber \\
        &\quad + \int_\Omega (\beta_1 + b(x))(v_1(x))^2 \, dx \nonumber \\
        &= 0.
    \end{align}
    
    Since $0 \leq c \lim_{N \to \infty} \|v_1^N - v_1\|_{L^2}^2 \leq 0$ and $c > 0$, it immediately follows by the squeeze theorem that:
    \begin{equation}
        \lim_{N \to \infty} \|v_1^N - v_1\|_{L^2}^2 = 0.
    \end{equation}
    Thus, $v_1^N \to v_1$ strongly in $L^2(\Omega)$. Because strong convergence preserves the standard $L^2$-norm, we finally conclude $\|v_1\|_{L^2} = \lim_{N \to \infty} \|v_1^N\|_{L^2} = 1$, confirming that $v_1 \in A_1$ is a valid, non-trivial normalized eigenfunction.
\end{proof}

\end{document}
